\section{可复现性说明}
实验分三阶段保存：首先导出物理参数、五组共同初始点及精确的MATLAB扰动样本；随后按固定预算运行三种纯搜索器；最后在MATLAB中对十五个固定候选作两档复核，并以复核通过的纯PSO参考重跑MPC。所有原始收敛历史、参数、计数、搜索值和复核残差随实验归档。

主实验目录为\texttt{results/comparison/}。机器可读清单记录算法、优化种子、扰动种子、积分步长及预算；CSV逐项记录全部十五次运行。论文数值宏、表格和图由归档数据生成，不手工转录。历史PSO加SQP结果单独保留，不参与本文表格或控制结果。MATLAB纯算法入口与Python编译评估入口均保留；不同随机数实现产生的轨迹不要求逐位相同，复现本文数值应使用本文记录的Python搜索入口及导出输入。

全局搜索每次评价17376条候选；最终两次ODE复核为额外固定验收开销，不进行局部搜索。搜索计时包含初始化评价及各轮更新，不包含JIT预热、最终复核及MPC。图中误差的显示下限不进入统计。
